\documentclass[11pt,a4paper]{article}

\usepackage[utf8]{inputenc}

\usepackage[T1]{fontenc}

\usepackage[polish]{babel}

\usepackage{lmodern}

\usepackage{geometry}

\geometry{a4paper,margin=27mm}

\usepackage{microtype}

\usepackage[hidelinks]{hyperref}

% Portfolio: MSC-AR-P29; source: prospective ARC-XR research scope.

\title{Predykcja subiektywnej niepewności decyzji z mikroruchów dłoni, spojrzenia i EEG}

\author{Mateusz Pomianek}

\date{Wersja 0.3.0, 21 listopada 2026}

\begin{document}

\maketitle

\begin{abstract}

Przedmiotem opracowania jest Predykcja subiektywnej niepewności decyzji z mikroruchów dłoni, spojrzenia i EEG. Czy ruch dłoni, okulografia i EEG pozwalają prognozować subiektywną pewność decyzji? Zestawione zostaną modele pojedynczych modalności i model łączący EEG, okulografię i kinematykę dłoni, z oceną na niezależnych uczestnikach. Głównym punktem końcowym jest błąd predykcji skali pewności oraz kalibracja prawdopodobieństw ocen. Opis ma charakter wstępny i nie przedstawia niezweryfikowanych wyników.

\end{abstract}

\section{Wprowadzenie}

Rozwój adaptacyjnych systemów XR stwarza możliwość badania interakcji człowieka z informacją przestrzenną, ale samo zastosowanie AR, VR lub modeli sztucznej inteligencji nie przesądza o przewadze poznawczej czy ergonomicznej. W rozpatrywanym zagadnieniu, Predykcja subiektywnej niepewności decyzji z mikroruchów dłoni, spojrzenia i EEG, kluczowe jest rozdzielenie właściwości systemu od zmienności zachowania użytkownika. Porównanie należy zaprojektować tak, aby mierzyć jednoznaczne zjawisko, a nie tylko ogólną ocenę technologii.

\section{Cel pracy i pytanie badawcze}

Czy ruch dłoni, okulografia i EEG pozwalają prognozować subiektywną pewność decyzji? Projekt wymaga krytycznego przeglądu aktualnego stanu badań oraz zdefiniowania nowości w stosunku do innych tematów portfolio XR.

\section{Zakres i odrębny wkład naukowy}
Modelowanie jawnie raportowanej pewności decyzji z kinematyki, spojrzenia i EEG; nie klasyfikacja obciążenia poznawczego P01 ani diagnoza psychologiczna. Predykcja ilościowa samoopisu na osobach spoza treningu, nie odczyt myśli ani ten sam konstrukt co obciążenie.
\section{Hipotezy badawcze}
\paragraph{H1.} Fuzja kinetyki dłoni, sygnału spojrzenia i EEG zmniejsza błąd predykcji deklarowanej pewności względem najsilniejszego predyktora jednomodalnego.
\paragraph{H2.} Zysk utrzymuje się przy kontroli poprawności decyzji i trudności zadania, nie tylko przez odgadywanie tego czy odpowiedź jest poprawna.
\paragraph{H3.} Kalibracja wyjściowego rozkładu pewności zachowuje poprawne pokrycie na niewidzianych uczestnikach i sesjach.
\section{Projekt badania i plan eksperymentów}
\subsection{A. Operacjonalizacja i kwalifikacja danych}
Modelowanie jawnie raportowanej pewności decyzji z kinematyki, spojrzenia i EEG; nie klasyfikacja obciążenia poznawczego P01 ani diagnoza psychologiczna. Jednostka wnioskowania: uczestnik × decyzja z odrębnym zgłoszeniem pewności. Przed badaniem należy określić źródła, wersje urządzeń, czas akwizycji, etykiety i niezależną referencję. Nie liczyć zależnych okien i klatek jako niezależnych prób.
\subsection{B. Eksperyment porównawczy}
Na stanowisku ekranowym zsynchronizować sygnał Tobii Eye Tracker 5, kinematykę i EEG; mieć niezależną deklarację pewności oraz stan zadania. Wynik pierwszorzędowy: różnica średniego błędu przewidywania skali pewności między modelem fuzji a najsilniejszym pojedynczym źródłem na niewidzianych osobach. Zamrozić losowanie warunków, progi, mianownik punktu końcowego i kryteria wyłączeń bez dostępu do oceny końcowej.
\subsection{C. Kontrole i analiza odporności}
Wykluczyć przeciek etykiet ze zbioru testowego; dokumentować artefakty EEG i braki sygnału, bez twierdzenia o zintegrowanym eye-trackingu Quest 3. Rozdzielać poprawność decyzji od jej subiektywnej pewności oraz przejść przez odrębny przegląd zgód biometrycznych. Utrata sygnału, abstencja, przerwanie zadania i dane brakujące są odrębnymi kategoriami wyniku.
\subsection{D. Walidacja na niezależnym materiale}
Walidacja musi zachować podział po uczestniku, sesji, scenie lub zadaniu. Bootstrap uczestnikowy MAE/RMSE w skali pewności i błąd kalibracji przewidywanych przedziałów. Raportować wyniki negatywne i rozbieżności z założeniami.
\section{Estymand, analiza statystyczna i decyzje}
różnica średniego błędu przewidywania skali pewności między modelem fuzji a najsilniejszym pojedynczym źródłem na niewidzianych osobach. Jednostka: uczestnik × decyzja z odrębnym zgłoszeniem pewności. Bootstrap uczestnikowy MAE/RMSE w skali pewności i błąd kalibracji przewidywanych przedziałów. H1 jest planowanym kontrastem głównym; H2 i H3 pozostają pomocnicze do zatwierdzenia hierarchii testów oraz kontroli wielokrotności. Wartości marginesów, mocy, minimalnego efektu istotnego praktycznie i liczebności mają wynikać z pilotażu, a nie z post hoc wyboru wyniku.
\section{Stan badań i granica nowości}
\begin{itemize}
\item EEG/VR time alignment and experimental validity must be explicitly documented before any multimodal inference claim. \cite{xrprior1}.
\item Eye tracking in immersive VR requires explicit validity and construct assumptions; device integration cannot be presumed. \cite{xrprior2}.
\end{itemize}
Predykcja ilościowa samoopisu na osobach spoza treningu, nie odczyt myśli ani ten sam konstrukt co obciążenie. Odrębność na tle opublikowanych metod i pozostałych artykułów wymaga osobnego audytu.
\section{Warunki bezpieczeństwa, sprzętu i zgodności}
To prospektywna koncepcja, nie zgoda na eksperyment, prerejestracja lub wynik empiryczny. Praca z uczestnikami, interwencje stresowe i pomiar biometryczny wymagają przeglądu etycznego, świadomej zgody i ochrony danych. Quest 3 w wersji standardowej nie posiada eye trackingu, a Tobii Eye Tracker 5 jest urządzeniem ekranowym; nie zastępować jednego stanowiska drugim.
\section{Materiał i metody}

Zestawione zostaną modele pojedynczych modalności i model łączący EEG, okulografię i kinematykę dłoni, z oceną na niezależnych uczestnikach. Wymagane jest zdefiniowanie jednostki analizy, zasad doboru przypadków, wersji scenariusza i źródeł pomiaru. Rozkład warunków oraz reguły wykluczenia prób należy określić przed pomiarem. W badaniach z udziałem ludzi niezbędna jest uprzednia właściwa ocena etyczna i świadoma zgoda. W testach wyłącznie symulacyjnych należy oddzielić prawdę symulacyjną od obserwacji empirycznej.

\subsection{Wynik główny i kryteria jakości}

Błąd predykcji skali pewności oraz kalibracja prawdopodobieństw ocen. Pomiar powinien zachowywać jawne jednostki, znaczniki czasu i reguły oceny poprawności. Przy modelach predykcyjnych stosować rozłączne partycje oceny; przy eksperymentach interwencyjnych oddzielić skuteczność predykcji od przyczynowego wpływu interfejsu.

\section{Plan analizy}

Walidacja rozdzielona według uczestników, analiza braków modalności i jakości sygnałów. W wynikach należy raportować wielkość efektu lub błąd estymacji z przedziałem niepewności, liczbą obserwacji i odsetkiem odrzuconych prób. Analizy dodatkowe wymagają jasnego oznaczenia jako eksploracyjne. Nie wybierać progów decyzyjnych ani modelu na podstawie obserwacji ze zbioru testowego.

\section{Wyniki}

W obecnej wersji nie raportuje się wyników pomiarów, tabel statystycznych ani parametrów wyznaczonych z nieudokumentowanych danych. To kierunek badawczy, a nie zatwierdzony protokół badania ani zgłoszony artykuł.

\section{Dyskusja}

Interpretacja przyszłych wyników musi uwzględnić alternatywne mechanizmy i zależność wniosków od jakości danych, warunków zadania oraz ograniczeń sprzętu. Wartość potencjalnego rozwiązania zależy nie tylko od trafności modelu, lecz także od jego opóźnienia, stabilności i wpływu na zachowanie człowieka. Wyniki testów syntetycznych nie mogą być automatycznie uogólniane na populację uczestników.

\section{Ograniczenia badania}

Predykcja pewności nie oznacza odczytu myśli, a sygnały ze stanowiska ekranowego nie mogą być traktowane jako zapis z Quest 3. Metodologiczny wynik dotyczący pewności decyzji, nie przeciążenia w MSC-AR-P01 ani proxy wymagań zadania w MSC-AFFECT-P05. Żadna część szkicu nie dowodzi przeprowadzenia zbierania danych, uzyskania zgód etycznych ani gotowości manuskryptu do zgłoszenia.

\section{Wnioski}

Przedstawiono problem, kontrast i sposób jego planowanej oceny. Rozstrzygnięcie postawionego pytania wymaga przygotowania protokołu, zweryfikowanych źródeł oraz poprawnie udokumentowanych wyników. Bez tych materiałów nie można sformułować konkluzji empirycznych.

\begin{thebibliography}{9}
\bibitem{xrprior1} The integration of virtual reality and EEG: A step-by-step guideline (2026), Methods guideline. \url{https://pmc.ncbi.nlm.nih.gov/articles/PMC12808493/}.
\bibitem{xrprior2} Validity of Eye-Tracking Measures in Immersive Virtual Reality: Using Cognitive Theories of Learning (2026), Educational Psychology Review. \url{https://doi.org/10.1007/s10648-026-10186-9}.
\end{thebibliography}
\end{document}

